% Options for packages loaded elsewhere
\PassOptionsToPackage{unicode}{hyperref}
\PassOptionsToPackage{hyphens}{url}
\documentclass[
]{article}
\usepackage{xcolor}
\usepackage{amsmath,amssymb}
\setcounter{secnumdepth}{-\maxdimen} % remove section numbering
\usepackage{iftex}
\ifPDFTeX
  \usepackage[T1]{fontenc}
  \usepackage[utf8]{inputenc}
  \usepackage{textcomp} % provide euro and other symbols
\else % if luatex or xetex
  \usepackage{unicode-math} % this also loads fontspec
  \defaultfontfeatures{Scale=MatchLowercase}
  \defaultfontfeatures[\rmfamily]{Ligatures=TeX,Scale=1}
\fi
\usepackage{lmodern}
\ifPDFTeX\else
  % xetex/luatex font selection
\fi
% Use upquote if available, for straight quotes in verbatim environments
\IfFileExists{upquote.sty}{\usepackage{upquote}}{}
\IfFileExists{microtype.sty}{% use microtype if available
  \usepackage[]{microtype}
  \UseMicrotypeSet[protrusion]{basicmath} % disable protrusion for tt fonts
}{}
\makeatletter
\@ifundefined{KOMAClassName}{% if non-KOMA class
  \IfFileExists{parskip.sty}{%
    \usepackage{parskip}
  }{% else
    \setlength{\parindent}{0pt}
    \setlength{\parskip}{6pt plus 2pt minus 1pt}}
}{% if KOMA class
  \KOMAoptions{parskip=half}}
\makeatother
\usepackage{color}
\usepackage{fancyvrb}
\newcommand{\VerbBar}{|}
\newcommand{\VERB}{\Verb[commandchars=\\\{\}]}
\DefineVerbatimEnvironment{Highlighting}{Verbatim}{commandchars=\\\{\}}
% Add ',fontsize=\small' for more characters per line
\newenvironment{Shaded}{}{}
\newcommand{\AlertTok}[1]{\textcolor[rgb]{1.00,0.00,0.00}{\textbf{#1}}}
\newcommand{\AnnotationTok}[1]{\textcolor[rgb]{0.38,0.63,0.69}{\textbf{\textit{#1}}}}
\newcommand{\AttributeTok}[1]{\textcolor[rgb]{0.49,0.56,0.16}{#1}}
\newcommand{\BaseNTok}[1]{\textcolor[rgb]{0.25,0.63,0.44}{#1}}
\newcommand{\BuiltInTok}[1]{\textcolor[rgb]{0.00,0.50,0.00}{#1}}
\newcommand{\CharTok}[1]{\textcolor[rgb]{0.25,0.44,0.63}{#1}}
\newcommand{\CommentTok}[1]{\textcolor[rgb]{0.38,0.63,0.69}{\textit{#1}}}
\newcommand{\CommentVarTok}[1]{\textcolor[rgb]{0.38,0.63,0.69}{\textbf{\textit{#1}}}}
\newcommand{\ConstantTok}[1]{\textcolor[rgb]{0.53,0.00,0.00}{#1}}
\newcommand{\ControlFlowTok}[1]{\textcolor[rgb]{0.00,0.44,0.13}{\textbf{#1}}}
\newcommand{\DataTypeTok}[1]{\textcolor[rgb]{0.56,0.13,0.00}{#1}}
\newcommand{\DecValTok}[1]{\textcolor[rgb]{0.25,0.63,0.44}{#1}}
\newcommand{\DocumentationTok}[1]{\textcolor[rgb]{0.73,0.13,0.13}{\textit{#1}}}
\newcommand{\ErrorTok}[1]{\textcolor[rgb]{1.00,0.00,0.00}{\textbf{#1}}}
\newcommand{\ExtensionTok}[1]{#1}
\newcommand{\FloatTok}[1]{\textcolor[rgb]{0.25,0.63,0.44}{#1}}
\newcommand{\FunctionTok}[1]{\textcolor[rgb]{0.02,0.16,0.49}{#1}}
\newcommand{\ImportTok}[1]{\textcolor[rgb]{0.00,0.50,0.00}{\textbf{#1}}}
\newcommand{\InformationTok}[1]{\textcolor[rgb]{0.38,0.63,0.69}{\textbf{\textit{#1}}}}
\newcommand{\KeywordTok}[1]{\textcolor[rgb]{0.00,0.44,0.13}{\textbf{#1}}}
\newcommand{\NormalTok}[1]{#1}
\newcommand{\OperatorTok}[1]{\textcolor[rgb]{0.40,0.40,0.40}{#1}}
\newcommand{\OtherTok}[1]{\textcolor[rgb]{0.00,0.44,0.13}{#1}}
\newcommand{\PreprocessorTok}[1]{\textcolor[rgb]{0.74,0.48,0.00}{#1}}
\newcommand{\RegionMarkerTok}[1]{#1}
\newcommand{\SpecialCharTok}[1]{\textcolor[rgb]{0.25,0.44,0.63}{#1}}
\newcommand{\SpecialStringTok}[1]{\textcolor[rgb]{0.73,0.40,0.53}{#1}}
\newcommand{\StringTok}[1]{\textcolor[rgb]{0.25,0.44,0.63}{#1}}
\newcommand{\VariableTok}[1]{\textcolor[rgb]{0.10,0.09,0.49}{#1}}
\newcommand{\VerbatimStringTok}[1]{\textcolor[rgb]{0.25,0.44,0.63}{#1}}
\newcommand{\WarningTok}[1]{\textcolor[rgb]{0.38,0.63,0.69}{\textbf{\textit{#1}}}}
\usepackage{longtable,booktabs,array}
\newcounter{none} % for unnumbered tables
\usepackage{calc} % for calculating minipage widths
% Correct order of tables after \paragraph or \subparagraph
\usepackage{etoolbox}
\makeatletter
\patchcmd\longtable{\par}{\if@noskipsec\mbox{}\fi\par}{}{}
\makeatother
% Allow footnotes in longtable head/foot
\IfFileExists{footnotehyper.sty}{\usepackage{footnotehyper}}{\usepackage{footnote}}
\makesavenoteenv{longtable}
\setlength{\emergencystretch}{3em} % prevent overfull lines
\providecommand{\tightlist}{%
  \setlength{\itemsep}{0pt}\setlength{\parskip}{0pt}}
\usepackage[]{natbib}
\bibliographystyle{plainnat}
\usepackage{bookmark}
\IfFileExists{xurl.sty}{\usepackage{xurl}}{} % add URL line breaks if available
\urlstyle{same}
\hypersetup{
  pdftitle={Descriptor-Teacher Guided Multi-to-Uni Transfer for Low-Resource Regression ADMET Prediction},
  pdfauthor={Author names to be added},
  hidelinks,
  pdfcreator={LaTeX via pandoc}}

\title{Descriptor-Teacher Guided Multi-to-Uni Transfer for Low-Resource
Regression ADMET Prediction}
\author{Author names to be added}
\date{}

\begin{document}
\maketitle

\subsection{Abstract}\label{abstract}

Low-resource ADMET prediction is a recurring challenge in early-stage
drug discovery, where labeled assays are limited but inexpensive
molecular knowledge such as physicochemical descriptors is readily
available. We study whether descriptor knowledge can be transferred into
an ECFP-only neural predictor, so that descriptors guide training
without being required at inference time. We develop Descriptor-Teacher
AdapterFusion, a lightweight ECFP4-based student model that combines an
adapter-based pseudo-descriptor representation with descriptor
reconstruction and teacher prediction distillation from an
ECFP4+descriptor random forest. Across five regression ADMET datasets,
three low-resource training ratios, and three random seeds, fixed
teacher distillation with \texttt{lambda\_distill\ =\ 1.0} improves
AdapterFusion in 11/15 dataset-ratio settings, with an average RMSE
change of -0.0656 relative to the undistilled AdapterFusion baseline.
However, descriptor-access random forest models remain stronger in most
settings, indicating that the proposed method partially transfers
descriptor knowledge but does not replace direct descriptor access. A
lambda sweep and adaptive-weighting ablation further show that simple
validation-based adaptive distillation is not sufficient. These results
support a controlled descriptor-guided multi-to-uni transfer perspective
for low-resource regression ADMET prediction.

\subsection{1. Introduction}\label{introduction}

Accurate ADMET prediction is important for early drug discovery because
absorption, distribution, metabolism, excretion, and toxicity properties
affect which candidate molecules can progress beyond screening
\citep{Vamathevan2019, Butler2018, Xiong2021ADMETlab}. In practice, many
ADMET prediction tasks are low-resource: labels are expensive to obtain,
measurements vary across assays, and each property may have a limited
number of reliable annotated compounds. This makes it useful to study
lightweight models and auxiliary molecular knowledge that can improve
generalization under limited supervision.

Molecular descriptors provide one such source of auxiliary knowledge.
RDKit physicochemical descriptors such as molecular weight, LogP, TPSA,
hydrogen bond counts, rotatable bonds, and ring counts are inexpensive
to compute and often correlate with ADMET behavior \citep{RDKit2026}. A
direct way to exploit them is to concatenate descriptors with molecular
fingerprints or use descriptor-based traditional models. However, this
creates a descriptor-access inference setting. In contrast, a
multi-to-uni transfer setting asks whether descriptor knowledge can be
used during training while the deployed student model uses only a single
modality, here ECFP4 fingerprints \citep{Rogers2010}.

This paper studies descriptor-guided multi-to-uni transfer for
low-resource ADMET prediction. We begin from an ECFP4 MLP baseline and
compare simple auxiliary descriptor prediction with a structured
AdapterFusion student that uses an ECFP-derived pseudo-descriptor
adapter and a learned fusion gate. We then add teacher prediction
distillation from an \texttt{ECFP4\_Desc\_RF} model, using a
descriptor-access random forest as a teacher for an ECFP-only neural
student.

Our central question is not whether a lightweight neural model can beat
all descriptor-access baselines. Instead, we ask a narrower and more
controlled question: can descriptor-access knowledge improve an
ECFP-only student in low-resource regression ADMET? The answer is
positive but bounded. Teacher distillation improves AdapterFusion in
most regression settings, but the descriptor-access teacher and
descriptor-based RF baselines remain strong.

The contributions are:

\begin{itemize}
\tightlist
\item
  We formulate a lightweight descriptor-guided multi-to-uni transfer
  setup for low-resource ADMET prediction.
\item
  We compare ECFP-only descriptor prediction, AdapterFusion,
  descriptor-access controls, random forest baselines, and
  teacher-distilled AdapterFusion under shared low-resource splits.
\item
  We show that fixed descriptor-teacher distillation improves regression
  AdapterFusion in 11/15 dataset-ratio settings.
\item
  We report lambda sensitivity, validation-selected lambda analysis, and
  a negative adaptive-weighting ablation, clarifying when simple teacher
  weighting is insufficient.
\end{itemize}

\subsection{2. Related Work}\label{related-work}

Molecular representation learning has increasingly used auxiliary
objectives to inject chemical knowledge into learned representations.
Molecular language models such as MolBERT and related approaches
motivate the use of domain-relevant pretraining or auxiliary tasks
rather than relying only on generic representation learning objectives
\citep{Fabian2020MolBERT}. Graph pretraining methods such as molecular
GNN pretraining, GROVER, and GraphMVP similarly use self-supervised or
cross-view objectives to improve downstream molecular prediction
\citep{Hu2020PretrainGNN, Rong2020GROVER, Liu2022GraphMVP}. Our work is
smaller in scale and does not use a large pretrained sequence encoder,
but it shares the principle that chemically meaningful auxiliary targets
can shape molecular representations.

Auxiliary learning and task-specific adaptation are also relevant. A
simple auxiliary objective may not transfer all useful knowledge to the
downstream task; the architecture through which auxiliary information is
routed can matter \citep{Ruder2017, Dey2024AuxAdapt}. In this work,
\texttt{DescPred} represents plain descriptor prediction, while
AdapterFusion represents a task-specific transfer mechanism that creates
and fuses a pseudo-descriptor representation.

The closest conceptual motivation is multi-to-uni molecular
representation learning, including M2UMol-style knowledge transfer
\citep{Xiong2026M2UMol}. Those works study how multimodal molecular
knowledge can improve unimodal inference. Our setting is deliberately
controlled: descriptors are the auxiliary modality, ECFP4 is the
inference modality, and ADMET prediction is the downstream task.

ADMET benchmark modeling provides the application context. We use
TDC-style ADMET datasets and evaluate low-resource splits across
regression and classification properties \citep{Huang2021TDC}.
MoleculeNet provides an earlier benchmarking reference point for
molecular property prediction metrics, splits, and baselines
\citep{Wu2018MoleculeNet}. The main method claim is limited to
regression ADMET because the observed distillation signal is more
consistent there than in classification.

Our baseline choices are also grounded in prior molecular property
prediction practice. Fingerprint and descriptor models remain important
controls in cheminformatics, while graph neural networks and directed
message-passing models provide stronger learned-representation baselines
in larger studies
\citep{Kearnes2016GraphConv, Yang2019DMPNN, Xiong2020AttentiveFP, Heid2024Chemprop}.
Tree ensembles are included because random forests and gradient boosting
methods remain competitive for tabular molecular features
\citep{Breiman2001, Chen2016XGBoost, Ke2017LightGBM, Tian2022ADMETBoost}.

\subsection{3. Method}\label{method}

\subsubsection{3.1 Problem Setting}\label{problem-setting}

Each molecule has an ECFP4 fingerprint \texttt{x\_ecfp}, optional
descriptor vector \texttt{x\_desc}, and target \texttt{y}. The primary
inference constraint is that the student model should use only
\texttt{x\_ecfp} at test time. Descriptor information may be used during
training as auxiliary supervision or through teacher predictions.

For regression tasks, the task loss is mean squared error. For
classification reference experiments, the task loss is binary cross
entropy. The main paper focuses on regression.

\subsubsection{3.2 ECFP Encoder}\label{ecfp-encoder}

The base ECFP model maps a 2048-dimensional ECFP4 fingerprint to a
128-dimensional latent representation using a two-layer MLP with batch
normalization, ReLU activations, and dropout:

\begin{Shaded}
\begin{Highlighting}[]
\NormalTok{z\_fp = f\_fp(x\_ecfp)}
\end{Highlighting}
\end{Shaded}

The \texttt{ECFP4\_MLP} baseline predicts the downstream property from
\texttt{z\_fp}.

\subsubsection{3.3 Descriptor Prediction
Baseline}\label{descriptor-prediction-baseline}

The \texttt{ECFP4\_MLP\_DescPred} baseline adds an auxiliary descriptor
prediction head to the ECFP encoder. The model predicts the downstream
target from \texttt{z\_fp} and predicts standardized RDKit descriptors
from the same representation.

The descriptor vector contains nine RDKit descriptors:

\begin{Shaded}
\begin{Highlighting}[]
\NormalTok{MolWt, LogP, TPSA, HBA, HBD, RotBonds, AromaticRings, HeavyAtoms, RingCount}
\end{Highlighting}
\end{Shaded}

The training objective is:

\begin{Shaded}
\begin{Highlighting}[]
\NormalTok{L = L\_task + lambda\_transfer * L\_desc}
\end{Highlighting}
\end{Shaded}

where \texttt{lambda\_transfer\ =\ 0.1}.

\subsubsection{3.4 AdapterFusion Student}\label{adapterfusion-student}

The AdapterFusion student uses the same ECFP encoder but introduces a
pseudo-descriptor adapter:

\begin{Shaded}
\begin{Highlighting}[]
\NormalTok{z\_desc = f\_adapter(z\_fp)}
\end{Highlighting}
\end{Shaded}

A learned fusion gate combines the original ECFP representation and the
pseudo-descriptor representation:

\begin{Shaded}
\begin{Highlighting}[]
\NormalTok{g = sigmoid(f\_gate([z\_fp, z\_desc]))}
\NormalTok{z\_fused = g * z\_fp + (1 {-} g) * z\_desc}
\end{Highlighting}
\end{Shaded}

The target prediction head receives \texttt{z\_fused}, while the
descriptor prediction head predicts descriptors from \texttt{z\_desc}.
This structure keeps inference ECFP-only, because no true descriptors
are consumed by the student at test time.

\subsubsection{3.5 Descriptor-Access
Controls}\label{descriptor-access-controls}

We include descriptor-access controls to estimate the value of direct
descriptor information. \texttt{ECFP4\_MLP\_DescConcat} concatenates
ECFP4 and standardized descriptors as neural input. We also train random
forest baselines with ECFP4-only, descriptor-only, and ECFP4+descriptor
inputs.

The \texttt{ECFP4\_Desc\_RF} baseline is used as the descriptor-access
teacher for distillation. This follows the broader teacher-student idea
of knowledge distillation, where a deployed student is trained to match
predictions from a stronger or richer teacher \citep{Hinton2015}.

\subsubsection{3.6 Descriptor-Teacher
Distillation}\label{descriptor-teacher-distillation}

The main method, Descriptor-Teacher AdapterFusion, keeps the
AdapterFusion student architecture and adds a teacher prediction loss.
For regression, the teacher loss is:

\begin{Shaded}
\begin{Highlighting}[]
\NormalTok{L\_teacher = MSE(y\_student, y\_teacher)}
\end{Highlighting}
\end{Shaded}

The final loss is:

\begin{Shaded}
\begin{Highlighting}[]
\NormalTok{L = L\_task + lambda\_transfer * L\_desc + lambda\_distill * L\_teacher}
\end{Highlighting}
\end{Shaded}

The main configuration uses:

\begin{Shaded}
\begin{Highlighting}[]
\NormalTok{lambda\_transfer = 0.1}
\NormalTok{lambda\_distill = 1.0}
\end{Highlighting}
\end{Shaded}

This value is selected from the fixed lambda sweep over
\texttt{\{0.01,\ 0.1,\ 0.3,\ 1.0\}}. The model architecture is unchanged
relative to AdapterFusion; only the training objective is augmented with
teacher supervision.

\subsection{4. Experiments}\label{experiments}

\subsubsection{4.1 Datasets}\label{datasets}

The main regression benchmark includes five ADMET datasets:

\begin{itemize}
\tightlist
\item
  \texttt{caco2\_wang}
\item
  \texttt{lipophilicity\_astrazeneca}
\item
  \texttt{solubility\_aqsoldb}
\item
  \texttt{vdss\_lombardo}
\item
  \texttt{ppbr\_az}
\end{itemize}

Supplementary classification experiments include:

\begin{itemize}
\tightlist
\item
  \texttt{bbb\_martins}
\item
  \texttt{hia\_hou}
\item
  \texttt{pgp\_broccatelli}
\item
  \texttt{bioavailability\_ma}
\item
  \texttt{herg}
\end{itemize}

The main claims are restricted to regression tasks.

\subsubsection{4.2 Low-Resource Protocol}\label{low-resource-protocol}

Each dataset is evaluated with training ratios of 10\%, 20\%, and 50\%.
Each setting is run with seeds 42, 123, and 3407. All models use the
same dataset-ratio-seed splits. Descriptor scalers are fit only on the
training split to avoid leakage.

\subsubsection{4.3 Baselines}\label{baselines}

We compare:

\begin{itemize}
\tightlist
\item
  \texttt{ECFP4\_MLP}
\item
  \texttt{ECFP4\_MLP\_DescPred}
\item
  \texttt{ECFP4\_MLP\_DescAdapterFusion}
\item
  \texttt{ECFP4\_MLP\_DescConcat}
\item
  \texttt{ECFP4\_RF}
\item
  \texttt{Desc\_RF}
\item
  \texttt{ECFP4\_Desc\_RF}
\item
  Descriptor-Teacher AdapterFusion
\end{itemize}

\texttt{ECFP4\_MLP\_DescConcat}, \texttt{Desc\_RF}, and
\texttt{ECFP4\_Desc\_RF} use descriptors at inference and should be
interpreted as descriptor-access controls rather than ECFP-only
students.

\subsubsection{4.4 Metrics}\label{metrics}

Regression tasks use MAE and RMSE, with RMSE as the primary metric.
Classification tasks use ROC-AUC and PR-AUC, with ROC-AUC reported in
the supplementary classification table. All tables report mean
plus/minus standard deviation over three seeds.

\subsubsection{4.5 Ablations}\label{ablations}

We run three ablation analyses:

\begin{itemize}
\tightlist
\item
  Fixed distillation lambda sweep:
  \texttt{lambda\_distill\ in\ \{0.01,\ 0.1,\ 0.3,\ 1.0\}}.
\item
  Validation-selected lambda analysis, where the lambda with lowest mean
  validation RMSE is selected per dataset-ratio setting without
  retraining.
\item
  Adaptive teacher weighting using validation teacher advantage:
\end{itemize}

\begin{Shaded}
\begin{Highlighting}[]
\NormalTok{lambda\_effective = lambda\_max * max((base\_valid\_rmse {-} teacher\_valid\_rmse) / base\_valid\_rmse, 0)}
\end{Highlighting}
\end{Shaded}

The adaptive rule is treated as a negative ablation.

\subsection{5. Results}\label{results}

\subsubsection{5.1 Main Regression
Results}\label{main-regression-results}

Table 1 is provided in
\texttt{paper\_tables/table1\_regression\_main\_rmse.md}. The main
pattern is that descriptor-access RF models remain strong, confirming
that descriptor information is predictive for ADMET properties. However,
among ECFP-only neural models, AdapterFusion is consistently stronger
than simple descriptor prediction on regression tasks.

Before teacher distillation, AdapterFusion improves over DescPred in all
15 regression dataset-ratio settings. This supports the hypothesis that
descriptor knowledge is better transferred through a structured
pseudo-descriptor adapter and fusion mechanism than through a plain
auxiliary descriptor head alone.

\subsubsection{5.2 Descriptor-Teacher
Distillation}\label{descriptor-teacher-distillation-1}

Teacher distillation further improves AdapterFusion. With
\texttt{lambda\_distill\ =\ 1.0}, distilled AdapterFusion improves over
base AdapterFusion in 11/15 regression settings. The mean RMSE delta
versus base AdapterFusion is -0.0656.

The gains are most visible in settings where the descriptor-access
teacher has a large advantage over the ECFP-only neural student. For
example, \texttt{ppbr\_az} benefits from fixed strong distillation
across all three training ratios. Nevertheless, the distilled student
remains below the descriptor-access RF teacher in most settings. The
result should therefore be interpreted as partial descriptor knowledge
transfer, not replacement of descriptor-access baselines.

\subsubsection{5.3 Lambda Sensitivity}\label{lambda-sensitivity}

The lambda ablation is provided in
\texttt{paper\_tables/table2\_lambda\_ablation\_rmse.md} and summarized
in \texttt{paper\_tables/table3\_lambda\_summary.md}.

{\def\LTcaptype{none} % do not increment counter
\begin{longtable}[]{@{}
  >{\raggedright\arraybackslash}p{(\linewidth - 6\tabcolsep) * \real{0.2308}}
  >{\raggedright\arraybackslash}p{(\linewidth - 6\tabcolsep) * \real{0.2308}}
  >{\raggedright\arraybackslash}p{(\linewidth - 6\tabcolsep) * \real{0.2308}}
  >{\raggedleft\arraybackslash}p{(\linewidth - 6\tabcolsep) * \real{0.3077}}@{}}
\toprule\noalign{}
\begin{minipage}[b]{\linewidth}\raggedright
lambda\_distill
\end{minipage} & \begin{minipage}[b]{\linewidth}\raggedright
complete settings
\end{minipage} & \begin{minipage}[b]{\linewidth}\raggedright
beats AdapterFusion
\end{minipage} & \begin{minipage}[b]{\linewidth}\raggedleft
mean delta vs AdapterFusion
\end{minipage} \\
\midrule\noalign{}
\endhead
\bottomrule\noalign{}
\endlastfoot
0.01 & 15/15 & 8/15 & 0.0075 \\
0.1 & 15/15 & 8/15 & -0.0221 \\
0.3 & 15/15 & 10/15 & -0.0046 \\
1.0 & 15/15 & 11/15 & -0.0656 \\
\end{longtable}
}

The strongest global fixed value is \texttt{lambda\_distill\ =\ 1.0}.
Smaller lambdas can be better for individual settings, but they do not
improve aggregate performance.

\subsubsection{5.4 Validation-Selected
Lambda}\label{validation-selected-lambda}

The validation-selected lambda analysis is shown in
\texttt{paper\_tables/table6\_validation\_selected\_lambda.md}.
Selecting lambda by mean validation RMSE beats fixed
\texttt{lambda\_distill\ =\ 1.0} in only 2/15 settings. It beats base
AdapterFusion in 12/15 settings.

This suggests that validation-selected lambda is useful as an analysis
but does not provide a stronger default than fixed
\texttt{lambda\_distill\ =\ 1.0} under the current setup.

\subsubsection{5.5 Adaptive Distillation Negative
Ablation}\label{adaptive-distillation-negative-ablation}

The adaptive teacher-weighting result is shown in
\texttt{paper\_tables/table4\_adaptive\_summary.md} and
\texttt{paper\_tables/table5\_adaptive\_vs\_fixed.md}. The adaptive rule
improves over base AdapterFusion in 9/15 settings, but beats fixed
\texttt{lambda\_distill\ =\ 1.0} in only 3/15 settings and beats the
best fixed lambda in only 1/15 settings. Its mean RMSE delta relative to
fixed \texttt{lambda\_distill\ =\ 1.0} is 0.0817, indicating worse
average performance.

This negative result shows that raw validation teacher advantage is not
a sufficient reliability signal for adaptive distillation.

\subsubsection{5.6 Classification
Reference}\label{classification-reference}

Classification results are provided in
\texttt{paper\_tables/tableS1\_classification\_main\_roc\_auc.md}. The
results are less stable than regression results, and the current
teacher-distillation claim is therefore limited to regression ADMET.

\subsection{6. Discussion}\label{discussion}

The experiments support a controlled descriptor-guided transfer story.
Descriptor information is valuable, and descriptor-access baselines
remain strong. The contribution is not to replace descriptor-access RF,
but to show that descriptor-access knowledge can partially improve an
ECFP-only neural student under low-resource regression ADMET settings.

AdapterFusion improves over plain descriptor prediction, suggesting that
the structure of auxiliary transfer matters. Teacher distillation then
provides an additional improvement by aligning the ECFP-only student
with predictions from a descriptor-access RF teacher.

The lambda sweep shows that stronger teacher supervision is generally
useful: \texttt{lambda\_distill\ =\ 1.0} is the best fixed global
setting. However, the dataset-specific best lambda varies, and neither
validation-selected lambda nor the tested adaptive rule clearly improves
over the fixed setting. The adaptive failure is informative:
dataset-level validation teacher advantage can underweight teacher
supervision on settings where strong distillation is beneficial.

Practically, these results recommend a cautious use of
descriptor-teacher distillation. It can improve an ECFP-only regression
student, but descriptor baselines should still be included and reported.

\subsection{7. Limitations}\label{limitations}

This work has several limitations. First, the main claim is restricted
to regression ADMET tasks. Classification results are included only as
supplementary evidence because improvements are less consistent. Second,
the student model is lightweight and ECFP-based; the method has not yet
been tested with large pretrained molecular encoders. Third,
descriptor-access RF remains stronger than the distilled student in most
settings. Finally, the adaptive teacher-weighting strategy tested here
is simple and negative; more principled sample-level or learned teacher
weighting may be needed.

\subsection{8. Conclusion}\label{conclusion}

We presented Descriptor-Teacher AdapterFusion, a lightweight
descriptor-guided multi-to-uni transfer method for low-resource
regression ADMET prediction. The method trains an ECFP-only
AdapterFusion student with descriptor reconstruction and teacher
prediction distillation from an ECFP+descriptor random forest. The main
fixed setting, \texttt{lambda\_distill\ =\ 1.0}, improves AdapterFusion
in 11/15 regression dataset-ratio settings. At the same time,
descriptor-access RF baselines remain stronger, so the appropriate
conclusion is controlled descriptor knowledge transfer rather than
overall ADMET state-of-the-art performance.

\bibliography{references}

\end{document}
